\newpage
\section*{РЕЗЮМЕ СТАРТАП-ПРОЕКТА}
\addcontentsline{toc}{section}{РЕЗЮМЕ СТАРТАП-ПРОЕКТА}


Название: «Онлайн-площадка для сдачи вещей в аренду». Проект представляет собой веб-приложение, предназначенное для автоматизации процессов аренды различных товаров между частными лицами или компаниями. Разработанная система должна цифровизировать такие задачи как просмотр каталога доступных вещей, бронирование товаров на определенный срок, расчет стоимости аренды, обработка заявок, а также взаимодействие между арендодателями и арендаторами. В ходе работы были проанализированы современные аналоги и сервисы совместного потребления (шеринг-платформы), и выявлены их отрицательные стороны:

\begin{enumerate}
	\item Интерфейс. Многие существующие платформы перегружены элементами, рекламой и второстепенными функциями, что затрудняет быстрый поиск нужного товара и оформление аренды для новых пользователей.
	
	\item Высокие комиссии и скрытые платежи за аренду товара или услуги. Ряд популярных сервисов взимают значительную комиссию как с арендодателей, так и с арендаторов, либо требуют оформления платной подписки для доступа к базовом возможностям, что снижает экономическую выгоду от использования платформы.
	
	\item Отсутствие гибкости в настройке условий аренды вещей или услуг. Во многих системах отсутствует возможность гибко настроить условия аренды для конкретного товара или услуги: посуточная/почасовая аренда, размер залога, стоимость доставки, что ограничивает владельцев вещей.
	
	\item Неудобная система коммуникации между арендатором и арендодателем. Встроенные чаты или формы связи часто работают с задержками, либо требуют перехода в сторонние мессенджеры, что создает неудобства и риски для безопасности сделки.
	
\end{enumerate}

Создаваемое приложение для аренды вещей позволит избежать вышеперечисленные недостатки без негативного влияния на процесс работы и скорость выполнения поставленных задач, предлагая сбалансированное решение с понятным интерфейсом и прозрачными условиями.

Актуальность проекта связана с ростом популярности модели совместного потребления (шеринг-экономики) в России и мире. Большое количество людей предпочитают арендовать вещи для разового или краткосрочного использования вместо их покупки, однако существующие площадки не всегда отвечают потребностям пользователей в простоте, надежности и доступности.

Целью является разработка полнофункционального веб-приложения для аренды вещей, обеспечивающего удобный поиск, бронирование и сдачу товаров в аренду.

Отличительные особенности разрабатываемого приложения по сравнению с уже существующими решениями:
\begin{enumerate}
	\item Интерфейс с продуманной структурой каталога и фильтрацией, позволяющей быстро подобрать вещь по категории, цене, сроку аренды.
	
	\item Прозрачная система формирования стоимости с автоматическим расчетом итоговой цены в зависимости от количества дней аренды, с возможностью учета скидок и залога.
	
	\item Реализация личного кабинета для арендодателей и арендаторов с разделением возможностей: управление объявлениями, просмотр истории бронирований, отслеживание статусов заказов.
	
	\item Интеграция с серверным API для выполнения основных операций (получение списка товаров, отправка заявок на аренду, добавление в корзину) без перезагрузки страницы, что обеспечивает высокую скорость работы и комфорт пользователя.
	
	\item Адаптивная верстка, обеспечивающая корректное отображение и удобство использования сайта на всех типах устройств: персональных компьютерах, планшетах и смартфонах.
\end{enumerate}

В разработанном приложении для аренды вещей реализованы такие функции как: просмотр каталога товаров с фильтрацией, регистрация и авторизация пользователей, добавление новых объявлений о сдаче вещей в аренду, оформление заказа с выбором дат аренды, корзина аренды, личный кабинет с историей операций, а также административная панель для модерации объявлений и управления пользователями.

Идея создания проекта появилась в сентябре 2025 года, активная разработка была начата в феврале 2026 года. С апреля 2026 года производилось тестирование системы путем моделирования различных пользовательских сценариев: размещение объявлений, поиск вещей, оформление бронирований. Проект планируется закончить к маю 2026 года. Также будет осуществлена регистрация программного продукта для защиты авторских прав и обеспечения исключительности прав.

Бизнес-цели и перспективы дальнейшего развития проекта:
\begin{enumerate}
	\item Достижение 2 500 успешных завершённых сделок (аренд) в течение первых 12 месяцев работы площадки.
	
	\item Обеспечение доли повторных аренд (пользователь арендует второй раз и более) на уровне не менее 30\% от общего числа активных клиентов к концу 1-го года.
	
	\item Снижение доли спорных сделок и возвратов залогов до менее 5\% от общего числа аренд путём внедрения системы верификации и страховки.
	
	\item Достижение среднего чека одной аренды на уровне не менее 1 500 рублей.
	
	\item Обеспечение времени отклика арендодателя на запрос арендатора не более 2 часов в 90\% случаев (для удержания пользователей).
\end{enumerate}